\DocumentMetadata{
  pdfversion=2.0,
  pdfstandard=ua-2,
  lang=en-US,
  tagging=on,
  tagging-setup={math/setup=mathml-SE,tabsorder=structure}
}
\documentclass[11pt,letterpaper]{article}
\usepackage[margin=0.68in,includefoot]{geometry}
\usepackage{newcomputermodern}
\usepackage{amsmath,amscd,mathtools}
\usepackage{tikz,adjustbox}
\providecommand{\square}{\mdlgwhtsquare}
\usepackage[hidelinks]{hyperref}
\hypersetup{
  pdftitle={Analysis First-Year Exam, August 2017, Part 2},
  pdfauthor={University of Florida Department of Mathematics},
  pdfsubject={Graduate mathematics examination},
  pdfdisplaydoctitle=true,
  bookmarksopen=true
}
\setcounter{secnumdepth}{0}
\setlength{\parindent}{0pt}
\setlength{\parskip}{0.28em}
\setlength{\emergencystretch}{3em}
\setlength{\leftmargini}{2.15em}
\setlength{\leftmarginii}{2.65em}
\setlength{\leftmarginiii}{3em}
\pagestyle{empty}
\raggedbottom
\allowdisplaybreaks
\NewTaggingSocketPlug{block/list/label}{examproblem}{%
  \tagstructbegin{tag=Lbl}\tagstructend%
  \tagstructbegin{tag=\LBody}%
  \tagstructbegin{tag=\ProblemHeadingTag,title-o={\ProblemHeadingText}}%
  \tagmcbegin{tag=\ProblemHeadingTag,actualtext={\ProblemHeadingText}}%
  #2%
  \tagmcend%
  \tagstructend%
}
\newenvironment{examproblems}[2]{%
  \def\ProblemHeadingTag{#2}%
  \AssignTaggingSocketPlug{block/list/label}{examproblem}%
  \begin{enumerate}%
  \setcounter{enumi}{#1}%
  \renewcommand{\labelenumi}{\textbf{\theenumi.}}%
  \setlength{\topsep}{0.35em}%
  \setlength{\partopsep}{0pt}%
  \setlength{\itemsep}{0.48em}%
  \setlength{\parsep}{0.18em}%
}{\end{enumerate}}
\newenvironment{examsubparts}{%
  \AssignTaggingSocketPlug{block/list/label}{default}%
  \begin{enumerate}%
  \setlength{\topsep}{0.18em}%
  \setlength{\partopsep}{0pt}%
  \setlength{\itemsep}{0.16em}%
  \setlength{\parsep}{0.1em}%
}{\end{enumerate}}
\newenvironment{examinstructions}{%
  \par\begingroup\itshape
}{%
  \par\endgroup\vspace{0.18em}
}
\makeatletter
\renewcommand\subsection{\@startsection{subsection}{2}{0pt}%
  {-1.25ex plus -.25ex minus -.1ex}%
  {0.55ex plus .1ex}%
  {\normalfont\large\bfseries}}
\renewcommand{\@maketitle}{%
  \newpage\null\vskip -1em%
  {\footnotesize\bfseries
    \href{https://www.ufl.edu/}{University of Florida}, \href{https://math.ufl.edu/}{Department of Mathematics}\par}%
  \vskip -0.5em%
  \tagstructbegin{tag=H1,title={Analysis First-Year Exam, August 2017, Part 2}}%
  \tagpdfparaOff%
  \tagmcbegin{tag=H1}%
    {\LARGE\bfseries\href{https://gma.math.ufl.edu/exams/analysis-first-year/}{Analysis First-Year Exam}, August 2017, Part 2\par}%
  \tagmcend%
  \tagpdfparaOn%
  \tagstructend%
  \vskip 0.25em%
  \tagmcbegin{artifact}%
    \hrule height 1pt%
  \tagmcend%
  \vskip 1em%
}
\makeatother
\title{Analysis First-Year Exam, August 2017, Part 2}
\author{}
\date{}
\begin{document}
\maketitle
\thispagestyle{empty}
\begin{examinstructions}
Answer FOUR questions in detail. State carefully any results used without proof.
\end{examinstructions}
\begin{examproblems}{0}{H2}
\def\ProblemHeadingText{Problem 1.}
\item
Let \(f:[0,1]\to[0,1)\) be continuous. Decide whether it follows that
\[
\int_0^1 \frac{1}{1-f(t)}\,dt=\sum_{n=0}^{\infty}\int_0^1 f(t)^n\,dt,
\]
 giving proof or counterexample as appropriate. [Note the intervals carefully.]
\def\ProblemHeadingText{Problem 2.}
\item
Let \(f:\mathbb{R}\to\mathbb{R}\) be nonconstant; for each positive integer~\(n\) and each real~\(t\) define \(f_n(t)=f(nt)\). Prove that there exists no \(\varepsilon>0\) such that \((f_n)_{n=1}^{\infty}\) is equicontinuous on the interval \((-\varepsilon,\varepsilon)\).
\def\ProblemHeadingText{Problem 3.}
\item
Let \(f:\mathbb{R}\to\mathbb{R}\) be continuous. Either
\begin{examsubparts}
\item[\textnormal{(i)}]
show that there exists a sequence of polynomials converging pointwise to~\(f\) on~\(\mathbb{R}\)
\item[\textnormal{(ii)}]
or show that there need not exist a sequence of polynomials converging uniformly to~\(f\) on~\(\mathbb{R}\).
\end{examsubparts}
\def\ProblemHeadingText{Problem 4.}
\item
Let \((f_n)_{n=1}^{\infty}\) be a sequence of measurable real-valued functions. Prove that each of the following sets is measurable:
\begin{examsubparts}
\item[\textnormal{(i)}]
\(A=\{\omega:f_n(\omega)\to\infty\text{ as }n\to\infty\}\);
\item[\textnormal{(ii)}]
\(B=\{\omega:f_n(\omega)\text{ is eventually irrational}\}\);
\item[\textnormal{(iii)}]
\(C=\{\omega:f_n(\omega)>0\text{ for infinitely many }n\}\).
\end{examsubparts}
\def\ProblemHeadingText{Problem 5.}
\item
Let \(f:[0,\infty)\to\mathbb{R}\) be locally Lebesgue integrable and assume that \(f(t)\to1\) as \(t\to\infty\). Prove that for each positive integer~\(n\) we may define
\[
a_n=\frac{1}{n}\int_0^{\infty}e^{-t/n}f(t)\,dt\in\mathbb{R}
\]
 and prove that \(a_n\to1\) as \(n\to\infty\).

\par
Suggestion: Say \(|f(t)-1|\leq\varepsilon\) whenever \(t\geq A\) and split the integral.
\end{examproblems}
\end{document}
